\chapter*{Resumo}
\addcontentsline{toc}{chapter}{Resumo}
O presente Projeto Integrador descreve o projeto, a construção e a validação experimental de um protótipo de fazenda urbana vertical inteligente, desenvolvido no âmbito da linha de pesquisa do LIADE (Laboratório de Pesquisa Aplicada e Desenvolvimento Eletrônico) sobre agricultura urbana e periurbana. O objetivo do sistema é automatizar o cultivo em ambientes internos por meio do monitoramento e controle autônomo das principais variáveis ambientais --temperatura, umidade do ar, umidade do substrato e intensidade luminosa-- sem exigir conhecimentos técnicos especializados por parte do usuário. Foi projetada e construída uma estrutura física modular de vários níveis para cultivo em substrato tradicional, utilizando rúcula e espinafre como culturas modelo, selecionadas por sua adaptabilidade agronômica a condições internas. Uma Raspberry Pi 4 foi selecionada como unidade de processamento central, responsável por receber as leituras de sensores digitais e analógicos e por acionar os atuadores por meio de estágios de relés de estado sólido e eletromecânicos: painéis de iluminação LED de espectro completo, uma rede de irrigação localizada com coleta e recirculação de água, e ventilação forçada. Peças sob medida, projetadas por manufatura aditiva, resolveram os desafios de impermeabilização, drenagem e montagem mecânica ao longo de toda a estrutura. Do lado do software, foi implementada em Python uma lógica de controle modular --baseada em critérios de limiar e histerese, filtragem de sinal e redundância de sensores-- para manter o microclima e a irrigação dentro dos parâmetros desejados a todo momento, enquanto um algoritmo de orçamento acumulado regula o fotoperíodo artificial em função da integral diária de luz (DLI) recebida por cada cultura. O protótipo também integra telemetria em nuvem e uma interface web remota, permitindo o monitoramento contínuo e a configuração do sistema de acordo com o paradigma da Internet das Coisas (IoT). Os ensaios experimentais realizados ao longo de vários meses validaram a capacidade do sistema de manter um microclima estável e favorável ao desenvolvimento das culturas, enquanto as medições de consumo de potência confirmaram que o protótipo opera dentro de uma faixa energeticamente eficiente, comparável à de um computador de mesa padrão.
\vspace{1cm}
\noindent\textbf{Área temática:} Eletrônica Digital e Sistemas de Controle
\vspace{0.5cm}
\noindent\textbf{Disciplinas:} Eletrônica Digital 3, Eletrônica Industrial, Eletrônica Analógica 1 e Sistemas de Controle 1
\vspace{0.5cm}
\noindent\textbf{Palavras-chave:} fazenda vertical, agricultura urbana, irrigação automatizada, monitoramento ambiental, controle embarcado, Internet das Coisas, cidades inteligentes
